\scanpage{70}
\mspage{69}
% E: text|upper left|B Recollements ne donnant pas lieu
B) Recollements \inserted{ne} donnant \uline{pas} lieu à une application naturelle $\mathcal X\longrightarrow X$, i.e. l'isom. $\sigma:\alpha\longrightarrow\beta$ est \uline{opposé} à celle obtenu par transp. par. \unsure{ou} par composition de celle-ci avec antipodisme --- ou encore, l'isom. de transition $\underline\omega(D)\longrightarrow\underline\omega(D')$ est \deleted{opposée à} celle du transport parallèle. Deux cas à envisager encore
% E: text|left middle|1 Cas alpha=a'
1) Cas $\alpha=a'$. C'est le cas où on repère une op. de générisation \inserted{sur $D'$} par l'arc qui fait face à l'arc de spécialisation correspondant sur $D$ --- \unsure{c'était un cas que pour une} \unsure{raison ou une autre} j'avais \unsure{écarté|écrit} (c'était vraiment un \unsure{bel imbroglio|beau méli-mélo}, rétrospectivement!).
% E: text|bottom left continuing upper right|2 Cas alpha=a
2) Cas $\alpha=a$ --- C'est le cas où on repère une op. de générisation par l'arc \inserted{de gén.} déduit de l'arc de spécialisation correspondant par transport parallèle. Je me suis amplement convaincu qu'il est impossible de ``\unsure{suivre|fermer}'' une p.r. en lui appliquant les itérés de l'opération \unsure{$\rho$|$\varphi$} de ``tourner autour d'un trou (un sommet\unsure{\ldots})''.
% E: text|right middle|Je crois me rappeler maintenant
Je crois me rappeler maintenant que j'avais \unsure{écarté|écrit} aussi le cas 1), car il donne \unsure{bel et bien} des circuits qui ont l'air idiots --- notamment en partant d'une p.r. proche d'une supersingulière $D_i$, on s'écarte de $D_i$ au lieu de tourner autour comme on devrait naturellement.
% E: deletion|bottom right|Là je viens de revérifier
\deleted{Là je viens de revérifier, cet ennui se présente aussi bien pour $B1)$ que pour $B2)$, \unsure{et ce} pour \deleted{B} $A1)$, et encore, je viens de voir, pour $A2$ --- visiblement, dans ces deux}
